% !TeX root = ../main.tex
\section{Thai}

% Everything on these two slides is typed as plain Thai. There is no wrapping
% macro anywhere -- that is the point of them. If any of it renders blank,
% the automatic fallback in preamble/lang-thai.tex has stopped working.

\begin{frame}{Thai text / ข้อความภาษาไทย}
  \begin{block}{Which font are you seeing?}
    Built with \textbf{XeLaTeX}, the Thai below is \hl{Sarabun} (TH Sarabun
    New) and the Latin stays Fira Sans.\\
    Built with \textbf{pdfLaTeX}, it falls back to \hl{Garuda}.
  \end{block}

  สวัสดีครับ นี่คือตัวอย่างข้อความภาษาไทยบนสไลด์ เพื่อทดสอบว่าฟอนต์แสดงผลถูกต้องหรือไม่

  Mixed in a sentence: the Thai word ตัวอย่าง means ``example''.

  Weight carries across the script boundary: \textbf{Results / ผลลัพธ์}
  and \textit{italic / ตัวเอียง}.

  % This slide is full: any taller and the last line runs into the reference
  % printed at the bottom. Beamer does not warn when that happens -- check by
  % eye after adding anything here.
  %
  % A citation inside Thai text. The reference at the bottom is Latin, so the
  % font has to switch back on the way out of the Thai run -- if the script
  % boundary handling in preamble/lang-thai.tex were wrong, this is where it
  % would show.
  งานวิจัยด้านการจัดวางตัวอักษรไทย \cite{tantisiriwat2020} อธิบายเรื่องนี้ไว้โดยละเอียด
\end{frame}

\begin{frame}{Thai list / รายการภาษาไทย}
  \begin{itemize}
    \item ข้อแรก --- first item
    \item ข้อที่สอง --- second item
    \item ข้อที่สาม --- third item
  \end{itemize}

  \vspace{0.3em}
  % Stacked marks: sara above + tone above (ที่ ปี้ กื๊อ), and vowels below
  % (ปรุง). These are the glyphs that collide if the line spacing is wrong.
  การจัดวางตัวอักษรไทยต้องรองรับสระบนและวรรณยุกต์ เช่น ที่ ปี้ กื๊อ ให้แสดงผลได้ถูกต้อง

  \vspace{0.3em}
  % A paragraph long enough to wrap several times. This is the real test:
  % Thai has no spaces between words, so it checks both that the lines break
  % at word boundaries (no hyphen mid-word) and that consecutive lines are
  % far enough apart for the marks above and below to clear each other.
  มาตรฐานข้อมูลเปิดระดับ 5 ดาว เป็นกรอบการประเมินคุณภาพข้อมูลที่ใช้กันอย่างแพร่หลาย
  โดยพิจารณาจากรูปแบบการเผยแพร่ ความสามารถในการเข้าถึง
  และการเชื่อมโยงข้อมูลระหว่างแหล่งต่าง ๆ
  ซึ่งหน่วยงานภาครัฐจำเป็นต้องปรับปรุงกระบวนการจัดเก็บและเผยแพร่ข้อมูลให้สอดคล้องกับเกณฑ์ดังกล่าว
  เพื่อให้ประชาชนสามารถนำข้อมูลไปใช้ประโยชน์ได้อย่างเต็มที่และมีประสิทธิภาพสูงสุด
\end{frame}
